\item Un cascarón esférico tiene 3.00 cm de radio interior y 7.00 cm de radio exterior. Está hecho de material con conductividad térmica $k = 0.800\text{ W/m}\cdot^\circ\text{C}$. El interior se mantiene a $5^\circ\text{C}$ de temperatura y el exterior a $40^\circ\text{C}$. Después de un intervalo de tiempo, el cascarón alcanza un estado estable en el que la temperatura en cada punto dentro de él se mantiene constante en el tiempo.
(a) Explique por qué la rapidez de transferencia de energía $P$ debe ser la misma a través de cada superficie esférica, de radio $r$, dentro del cascarón y debe satisfacer
$$ \frac{dT}{dr} = \frac{P}{4\pi k r^2} $$
(b) A continuación, demuestre que
$$ \int_{5}^{40} dT = \frac{P}{4\pi k} \int_{0.03}^{0.07} r^{-2} dr $$
donde $T$ está en grados Celsius y $r$ está en metros. 
(c) Encuentre la rapidez de transferencia de energía a través del cascarón. 
(d) Demuestre que
$$ \int_{5}^{T} dT = 1.84 \int_{0.03}^{r} r^{-2} dr $$
donde $T$ está en grados Celsius y $r$ está en metros. 
(e) Obtenga la temperatura dentro del cascarón como una función del radio. 
(f) Encuentre la temperatura en $r = 5.00\text{ cm}$.
